\section{Introduction}\label{sec:introduction}

Let $A$ be a finite-dimensional algebra over a field $k$.
A finite-dimensional $A$-module $M$ is \emph{self-orthogonal} if
$\Ext_A^i(M,M)=0$ for every integer $i>0$.
Tachikawa's second conjecture asserts that every finite-dimensional
self-orthogonal module over a finite-dimensional self-injective
algebra is projective.  This question belongs to the homological
conjectures surrounding quasi-Frobenius algebras and dominant
dimension; see Tachikawa's 1973 monograph~\cite[pp.~VI--VII]{Tachikawa1973} and the account
of the conjecture in~\cite[p.~7288]{ChenXi2025TAMS}.
The Auslander--Reiten conjecture, originating in
\cite{AuslanderReiten1975}, asks for finite-dimensional algebras
whether the condition $\Ext_A^i(M,M\oplus A)=0$ for every $i>0$
forces $M$ to be projective. For a self-injective algebra, the groups
$\Ext_A^i(M,A)$ vanish automatically, so Tachikawa's second conjecture
is its self-injective case.

Write $D=\Hom_k(-,k)$. An algebra $A$ is \emph{symmetric} if
$A\simeq DA$ as $A$-bimodules; in particular, a symmetric algebra
is self-injective. We prove the following result in this class.

\begin{theorem}\label{thm:main}
Let
\[
 k=\mathbb F_2(q,H_1,H_2),
\]
where $q,H_1,H_2$ are algebraically independent.
There exist a finite-dimensional associative unital symmetric
$k$-algebra $A$ and a finite-dimensional unital left $A$-module $M$
such that $M$ is nonprojective and
\[
 \Ext_A^i(M,M)=0\qquad\text{for every integer }i>0.
\]
In particular, Tachikawa's second conjecture is false.
\end{theorem}

Since $A$ is self-injective, the same module satisfies
$\Ext_A^i(M,M\oplus A)=0$ for every $i>0$.
Thus Theorem~\ref{thm:main} also gives a symmetric counterexample to
the Auslander--Reiten conjecture.

\paragraph{Related results.}
Schulz's earlier example of a nonprojective module without positive
self-extensions uses an algebra over a skew field, outside the
finite-dimensional-over-a-field setting considered here; see
\cite{Schulz1994} and the distinction recalled in
\cite[Introduction]{DiverisPurin2012}.
His quantum-exterior example has nonzero first self-extension but no
higher self-extensions
\cite[Example~7]{Schulz1995}; the parameter-shifting syzygies in that
example are an antecedent for our starting algebra.
For symmetric Artin algebras, Diveris and Purin show that a component
containing a nonprojective module with eventually vanishing
self-extensions has stable part of type $\mathbb Z A_\infty$
\cite[Theorem~A]{DiverisPurin2014}.
Over an algebraically closed field, Erdmann classifies, up to Morita
equivalence, the indecomposable
weakly symmetric algebras with radical cube zero and nonzero radical
square that admit such modules. The classification singles out quantum
exterior and double Nakayama algebras with parameters that are not roots
of unity~\cite[Theorem~1.1]{Erdmann2017}.

Chen, Fang and Xi characterize the symmetric case through stratifying
ideals and give conditional mirror-reflective constructions from
gendo-symmetric counterexamples to the Nakayama conjecture
\cite[Theorems~1.1 and~1.2]{ChenFangXi2024}.
Chen and Xi relate self-orthogonal Nakayama-stable generators over
self-injective algebras to additional categorical and
Gorenstein--Morita conditions~\cite[Theorem~1.2 and Corollary~1.4]{ChenXi2025TAMS}.
Their further criterion assumes virtual Gorensteinness of the relevant
endomorphism algebra~\cite[Theorem~1.3]{ChenXi2025VG}.
Zhang and Zhou prove the self-injective case for split algebras with
radical cube zero~\cite[Corollary~4.2]{ZhangZhou2026} and subsequently
treat Igusa--Todorov algebras~\cite[Theorem~3.2]{ZhangZhouIgusaTodorov2026}.
Xia proves the Auslander--Reiten conjecture for quantum complete
intersections~\cite[Corollary~4.3]{Xia2026}.
In a different direction, B\"ohmler and Marczinzik construct a
nonprojective module over a commutative local self-injective algebra
whose first two self-extension groups vanish, while the third does
not~\cite[Theorem~3.1 and the end of Section~3]{BohmlerMarczinzik2026}.
Theorem~\ref{thm:main} establishes vanishing in every positive degree.

Enomoto and Marczinzik prove that Tachikawa's second conjecture for the
two-fold trivial extension of an Artin algebra implies the Auslander--Reiten
conjecture for that algebra~\cite[Theorem~3.5]{Enomoto2026}.
This reduction has a self-injective target, not a symmetric target in
general. Our construction obtains a symmetric algebra by establishing
the additional complete-resolution vanishing needed for transfer to
the ordinary trivial extension.
Enomoto and Marczinzik also prove an equivalence among the universal
versions of the Auslander--Reiten, generalized Nakayama, Nakayama,
and both Tachikawa conjectures
over a fixed commutative Artinian coefficient
ring~\cite[Theorem~1.1]{Enomoto2026}.

\paragraph{The construction.}
The companion article~\cite{OpenAIAR2026} constructs an explicit
eight-simple triangular counterexample. We use the same ten-dimensional
starting algebra and related cone and triangular-module constructions,
and reproduce the algebra, resolutions, and arguments needed here.
The passage to a symmetric algebra requires vanishing in negative
degrees of a complete self-extension complex as well as positive Ext
vanishing.

Two steps organize the proof.  First, Proposition~\ref{tr:transfer}
gives a general transfer principle. For a finite-dimensional algebra
$\Lambda$, a doubly infinite exact complex $P_Z$ of finite projective
$\Lambda$-modules is \emph{totally acyclic}
if $\Hom_\Lambda(P_Z,\Lambda)$ is also exact. Suppose
$P_Z$ is totally acyclic, $Z=\coker((P_Z)_1\to(P_Z)_0)$, and
\begin{equation}\label{intro:profile}
 H^a\Hom_\Lambda(P_Z,Z)=0
 \qquad(a>0\ \text{or}\ a\leq-2).
\end{equation}
Then the trivial extension and induced module
\[
 A=\Lambda\ltimes D\Lambda,
 \qquad M=A\otimes_\Lambda Z
\]
give the required symmetric example whenever $Z$ is nonprojective.
The negative tail computes the positive extensions into
$D\Lambda\otimes_\Lambda Z$, with the precise index $-a-1$.
This transfer principle applies to any finite-dimensional algebra
and module satisfying its hypotheses.

Second, we construct the data required by this transfer. A
ten-dimensional algebra $C$ has a one-dimensional simple module $s$.
Put
\[
 T=C\ltimes DC,\qquad B=C\otimes_k C,
 \qquad E=T\otimes_k T,\qquad X=s\otimes_k s.
\]
We regard $s$ as a $T$-module through $T\to C$, and $X$ as an
$E$-module. Write $s[r]$ for $s$ placed in homological degree $r$.
Derived induction of $s$ from $C$ to $T$ gives an object
$\mathcal J=T\otimes_C^{\mathbf L}s$ in a triangle
\[
 s[2]\longrightarrow\mathcal J\longrightarrow s
 \xrightarrow{\tau}s[3].
\]
The class $\tau$ generates $\Ext_T^*(s,s)=k[\tau]$.
The tensor product $E$ is symmetric, and external products give the
self-extension algebra $\Ext_E^*(X,X)=k[\tau_1,\tau_2]$, where
$\tau_1=\tau\otimes1$ and $\tau_2=1\otimes\tau$ both have degree three.
A stable map is a
module map modulo maps factoring through projectives; over a symmetric
algebra the shift $[1]$ is inverse syzygy. Symmetric duality determines
the negative stable self-maps of $X$ from this polynomial algebra.

Tensoring two copies of this triangle describes
$E\otimes_B^{\mathbf L}X\simeq\mathcal J\otimes\mathcal J$.
Its attaching maps act by $\tau_1$ and $\tau_2$. On nonnegative
stable self-maps these give the two-variable Koszul sequence, and on
negative self-maps they give its dual. This Koszul calculation, together
with the full attaching maps, leaves only two scalar groups. The bar
complex for $E\otimes_B^{\mathbf L}E$ provides a finite $E$-bimodule
$Y$, projective separately on the left and right, whose evaluation
represents $(\mathcal J\otimes\mathcal J)[-1]$ in the stable category.
Its stable Hom groups are
\[
 W^a:=\sHom_E(X,(Y\otimes_E X)[a])\simeq
 \begin{cases}k,&a=-3,0,\\0,&a\notin\{-3,0\}.\end{cases}
\]
For $H\in k^\times$, let $h_H$ fix $C$ and multiply $DC$ by $H$,
and let $U_H$ be $E$ with right action twisted by $h_H\otimes h_H$.
There is a natural identification $U_H\otimes_E X=X$.
To combine two twists, we need the nonzero line $W^0$ to be reached
by evaluating genuine bimodule maps $U_H\to Y$; a module map
$X\to Y\otimes_E X$ alone does not supply a bimodule fiber.
A trace obstruction for a character map $B\to DB$, combined with
symmetric duality, produces precisely these maps $g_H$, with nonzero
evaluations in degree zero. The sparse Hom profile and these bimodule
lifts are the two conclusions of Proposition~\ref{b:bimodule}.

After normalizing the maps for $H_1,H_2$, their fiber has a finite
bimodule representative $F$, obtained as an ordinary kernel after
adjoining a projective summand. It comes with a stable map
$v:X\to F\otimes_E X$. The two twists act by distinct scalars on each
nonzero positive or negative homogeneous self-map space of $X$.
These actions make the self-extension comparison invertible in the
required degree ranges. Forming a cone over
\[
 \Lambda=\begin{pmatrix}E&0\\ F&E\end{pmatrix}
\]
produces the totally acyclic resolution of $Z$ with
\eqref{intro:profile}; the boundary degrees $1$ and $-2$ are checked
separately. Figure~\ref{fig:proof-architecture} locates the bimodule
construction and the final transfer in this argument.
All bimodule representatives, kernels, and cokernels used to define
the final algebra and module are finite-dimensional over the stated
field.  Infinite complexes serve to prove their properties.

\begin{figure}[ht]
\centering
\begin{tikzpicture}[>=Stealth,
 box/.style={draw,rounded corners=2pt,align=center,text width=4.8cm,
             inner sep=6pt,font=\small},
 lab/.style={font=\footnotesize,align=center}]
\node[box] (ex) at (0,3.7)
 {$E=T\otimes T$, $X=s\otimes s$\\$\Ext_E^*(X,X)=k[\tau_1,\tau_2]$};
\node[box] (y) at (6.6,3.7)
 {$Y$ and $g_H:U_H\to Y$\\$W^a=0$ for $a\notin\{-3,0\}$};
\node[box] (fv) at (6.6,1)
 {$F$ and $v:X\to F\otimes_E X$};
\node[box] (lz) at (0,1)
 {$(\Lambda,Z)$ and $P_Z$\\$H^a\Hom_\Lambda(P_Z,Z)=0$\\for $a>0$ or $a\leq-2$};
\node[box] (am) at (0,-1.6)
 {$A=\Lambda\ltimes D\Lambda$\\$M=A\otimes_\Lambda Z$\\symmetric counterexample};
\draw[->] (ex)--node[above=9mm,lab]{derived induction\\and trace obstruction}(y);
\draw[->] (y)--node[right,lab]{two twists\\and finite fiber}(fv);
\draw[->] (fv)--node[above=9mm,lab]{triangular cone}(lz);
\draw[->] (lz)--node[right,lab]{trivial extension\\and induction}(am);
\end{tikzpicture}
\caption{Constructions and implications in the proof. The complete
self-Hom vanishing for $(\Lambda,Z)$ supplies the hypotheses for
transfer to $(A,M)$. Arrows indicate proof stages, not algebra maps.}
\label{fig:proof-architecture}
\end{figure}

\paragraph{Homological consequences.}
The pair $(A,M)$ gives counterexamples to several conjectures formulated
through injective resolutions and tilting modules. Corollary~\ref{cor:homological}
realizes five such failures using one endomorphism algebra and two
associated modules. We state the relevant conjectures in their
finite-dimensional left-module forms before giving the corollary.
For a finite-dimensional algebra $\Gamma$, write
\[
 0\longrightarrow {}_\Gamma\Gamma\longrightarrow I^0
 \longrightarrow I^1\longrightarrow\cdots
\]
for its minimal injective resolution. Its \emph{dominant dimension} is
the supremum of the integers $d\geq0$ for which $I^j$ is projective
whenever $0\leq j<d$, and is infinite when every $I^j$
is projective. The classical \emph{Nakayama conjecture} asserts that
infinite dominant dimension forces self-injectivity.
The \emph{generalized Nakayama conjecture} asserts that every
indecomposable injective occurs as a summand of some $I^j$; equivalently,
for every simple module $S$, some $\Ext_\Gamma^i(S,\Gamma)$ with
$i\geq0$ is nonzero. The \emph{strong Nakayama conjecture} asserts
that a finite module $V$ with $\Ext_\Gamma^i(V,\Gamma)=0$ for all
$i\geq0$ is zero. The \emph{Auslander--Gorenstein conjecture} asserts
that the bounds $\pd_\Gamma I^j\leq j$ for every $j\geq0$ force
finite self-injective dimension on both sides. These formulations are
recalled in~\cite[Introduction]{Zhang2018HomologicalConjectures}.

Write $\add W$ for the direct summands of finite direct sums of a
module $W$. A finite left $\Gamma$-module $W$ is \emph{tilting} if it
has finite projective dimension, is self-orthogonal, and admits a
finite exact coresolution
\[
 0\longrightarrow\Gamma\longrightarrow W^0\longrightarrow\cdots
 \longrightarrow W^r\longrightarrow0,\qquad W^j\in\add W.
\]
It is \emph{Wakamatsu tilting} if it is self-orthogonal and admits an
exact coresolution
\[
 0\longrightarrow\Gamma\longrightarrow W^0\longrightarrow W^1
 \longrightarrow\cdots,\qquad W^j\in\add W,
\]
which remains exact after applying $\Hom_\Gamma(-,W)$. This is the
finite left-module definition in~\cite[Definition~2.1]{Huang2006}.
The \emph{Wakamatsu tilting conjecture} asserts that a Wakamatsu tilting
module of finite projective dimension is tilting; see
\cite{DivaaniAazarMahinFallahTousi2025} for a recent discussion.

\begin{corollary}[Homological conjectures]\label{cor:homological}
Let $A,M$ be as in Theorem~\ref{thm:main}. For every field extension
$K/k$, put $A_K=K\otimes_k A$, $M_K=K\otimes_k M$, and
\[
 \Gamma_K=\End_{A_K}(A_K\oplus M_K)^{\mathrm{op}}.
\]
Then $\Gamma_K$ is a finite-dimensional $K$-algebra with the following
properties, in the category of finite left $\Gamma_K$-modules.
\begin{enumerate}
\item Every term of the minimal injective resolution of $\Gamma_K$
is projective, but $\idim_{\Gamma_K}\Gamma_K=\infty$.
\item There is a nonzero simple module $S$ such that
$\Ext_{\Gamma_K}^i(S,\Gamma_K)=0$ for every $i\geq0$.
\item There is a projective-injective Wakamatsu tilting module $W$
which is not tilting.
\end{enumerate}
Consequently the classical, generalized, and strong Nakayama conjectures,
the Auslander--Gorenstein conjecture, and the Wakamatsu tilting conjecture
are false for finite-dimensional algebras over $K$.
\end{corollary}

The opposite in this definition keeps our left-module convention. The
endomorphism-algebra connection with dominant dimension is classical,
going back to Mueller~\cite[Lemma~3 and Theorem~4]{Mueller1968}; the precise
Tachikawa--Nakayama correspondence for a pair $(A,M)$ is also stated
in~\cite[p.~7288]{ChenXi2025TAMS}. Section~\ref{sec:consequences}
uses this connection and constructs the modules $S$ and $W$ directly.
Taking $K$ to be an algebraic closure of $k$ gives examples over an
algebraically closed field of characteristic two. This field contains
the displayed transcendental parameters; the assertion does not concern
the algebraic closure of a finite field.

The same resolution also gives a finitistic-dimension consequence.
For a finite-dimensional algebra, the \emph{little left finitistic
dimension} is the supremum of the finite projective dimensions of
finite left modules; the \emph{big left finitistic dimension} allows
all left modules instead.

\begin{corollary}[Finitistic dimensions]\label{cor:finitistic}
Let $k$ be the field in Theorem~\ref{thm:main}, let $K/k$ be any field
extension, and let $\Gamma_K$ be the algebra in
Corollary~\ref{cor:homological}. For every integer $n\geq1$, there is a
finite-dimensional left $\Gamma_K$-module $C_n$ such that
\[
 \pd_{\Gamma_K}C_n=n.
\]
Consequently $\Gamma_K$ has infinite little left finitistic dimension
and infinite big left finitistic dimension.
\end{corollary}

This gives an independent characteristic-two example alongside the
complex algebra of~\cite[Theorem~1.1]{OpenAIFinitistic2026}.

Two further applications are developed in
Section~\ref{sec:consequences}. The connection between generalized
Nakayama failure and cotilting complements established by Buan and Solberg
\cite[Theorems~3.1 and~3.4]{BuanSolberg1998} gives a projective module
over $\Gamma_K^{\mathrm{op}}$ that can be completed to a tilting module
in infinitely many ways, each by adjoining one indecomposable summand;
these summands are pairwise nonisomorphic
(Corollary~\ref{cor:tilting-complements}). Their introduction credits
an independent proof of this connection to Happel and Unger
\cite[Introduction]{BuanSolberg1998}. The mirror-reflective
constructions of Chen, Fang and Xi~\cite{ChenFangXi2024} give auxiliary
indecomposable symmetric algebras with arbitrarily long idempotent
stratifications and the resulting chains of successive derived-category
recollements specified in Corollary~\ref{cor:stratifying}.

\paragraph{Organization.}
Section~\ref{sec:homological} establishes the complete-resolution,
duality, and finite-representative tools.
Section~\ref{sec:transfer} proves the transfer to a symmetric algebra.
Sections~\ref{a:section} and~\ref{e:section} construct the starting
algebra and compute the stable self-maps and twist actions.
Section~\ref{b:section} constructs the finite bimodule and proves the
trace obstruction. Section~\ref{sec:triangular} forms the triangular
module, establishes the complete vanishing profile, and applies the
transfer to prove Theorem~\ref{thm:main}.
Section~\ref{sec:consequences} studies the associated endomorphism
algebra and proves Corollaries~\ref{cor:homological} and~\ref{cor:finitistic}.

\paragraph{Conventions.}
All algebras are associative and unital, all modules are unital left
modules unless indicated otherwise, and all modules occurring in the
construction are finite-dimensional.
Unspecified tensor products are over $k$, and
$G^e=G\otimes_k G^{\mathrm{op}}$ acts on $G$-bimodules by the usual
outer actions.  Complexes are indexed homologically, with
$P[a]_j=P_{j-a}$.  Since the field has characteristic two, signs in
their differentials and comparison maps disappear.  The following
section fixes the stable and derived conventions used throughout.
